\chapter{Diskussion}
Die Ergebnisse zeigen für die konkret gewählten Konfigurationen einen deutlichen Genauigkeitsvorteil von POD--GPR. Dies ist keine allgemeine Überlegenheit von Gaußprozessen: Die Aussage betrifft einen begrenzten Parameterraum, ein linear-elastisches Referenzproblem und die durchsuchten Modell- und Hyperparameterräume. Das direkte MLP wurde nur für die Breiten 64 und 128 untersucht; die ursprünglich vorgesehenen Breiten 32 und 80 fehlen. Das Ergebnis der RBF-Familie betrifft feste Gaußzentren und nachgeschaltete Tanh-Schichten.

Im kontrollierten Vollpoolvergleich ist der Projektionsfehler bereits bei Rang 2 klein. Dass der vollständige POD--GPR-Fehler bei höheren Rängen nur geringfügig sinkt, legt einen günstigen Kompromiss bei kleinem Rang nahe. Diese Interpretation entstand nach Einsicht in die Testdaten. Sie ist eine Folgehypothese und ändert nicht den auf der Validierung gewählten Rang 15. Bei POD--MLP verändern sich die Fehler je Rang anders; aus einzelnen kontrollierten Läufen sollte nicht auf die gesamte Verteilung möglicher neuronaler Initialisierungen geschlossen werden.

Die Extrapolationsfälle überschreiten jeweils nur eine Grenze. Die großen relativen Fehler bei kleiner Referenznorm sind nicht gleichbedeutend mit gleich großen absoluten Fehlern. Umgekehrt zeigen die absoluten Fehler und falschen Vorzeichen im auffälligen Lastfall, dass die relativen Ausreißer nicht allein ein Effekt der Normierung sind. Die nachträgliche Gruppenauswertung enthält nur zehn oder elf Fälle pro Gruppe und erlaubt keine gesicherte Ursachenzuordnung.

Die FEM-Felder sind Referenzen des verwendeten numerischen Netzes. Die Netzuntersuchung im Pilot ist begrenzt; eine flächendeckende Konvergenzprüfung liegt nicht vor. Analytische Balkentheorien dienen als Plausibilisierung und beschreiben nicht alle Komponenten des dreidimensionalen Feldes exakt. Die unabhängigen Testfälle sind in ihrer Rolle von der Validierung getrennt, wurden jedoch über dieselben Modelle wiederholt vorhergesagt; Seeds erzeugen keine zusätzlichen unabhängigen Balkenfälle.

Die technische Auswertekorrektur nach dem ersten Testzugriff ist eine Einschränkung der ursprünglichen Vorabprovenienz. Ihr Änderungsumfang, die unveränderten Modelle und Kennzahlen sowie der Vergleich der ersten beiden Split-Ergebnisse sind separat dokumentiert. Für POD--GPR wurde keine Unsicherheitskalibrierung vorgenommen; der Spitzen-RAM wurde nicht gemessen. Die Laufzeit- und Break-even-Werte hängen von Messumfang, Rechner, Datenkosten und Aufgabenverteilung ab. Der kleine spätere Benchmark erlaubt eine bedingte Szenariorechnung, nicht die Behauptung einer universellen Schwelle.

\chapter{Fazit und Ausblick}
Für den vorliegenden statischen Vergleich ist POD--GPR bei $N=768$ der genauigkeitsorientierte Ausgangspunkt: Der mittlere Interpolationsfehler beträgt 0,55\,\% und der mittlere Extrapolationsfehler 3,85\,\%. Dem steht eine vollständige Fit-Zeit des ausgewerteten Hauptmodells von 353,57\,s gegenüber. POD--MLP ist schneller trainiert, erreicht aber einen größeren Feldfehler. Direkte MLP- und RBF-Modelle sind weitere Bezugspunkte, deren Ergebnisse von ihrem begrenzten Suchraum abhängen. Der Datenbedarf wird nicht allein durch die Fallzahl bestimmt: Die separate Verteilungsstudie zeigt für alle vier Familien einen Vorteil des konkret betrachteten raumfüllenden Designs.

Für weiterführende Arbeiten bietet sich eine vorab definierte Untersuchung kleiner POD-Ränge mit neuen unabhängigen Bewertungsfällen an. Ebenso können Netzkonvergenz, kritische Lastfälle und eine mögliche Kalibrierung der GPR-Unsicherheit vertieft werden. Eine Erweiterung auf Eigenfrequenzen wäre ein eigenständiger nächster Schritt und ist nicht durch die statischen Ergebnisse dieser Arbeit belegt.
